\documentclass[11pt]{article}
\usepackage[landscape,margin=0.8in]{geometry}
\usepackage{amsmath,amssymb}
\usepackage{graphicx}
\usepackage{booktabs}
\usepackage{caption}
\usepackage{float}
\usepackage{xcolor}
% The auto-generated \input tables still carry some Japanese header labels,
% so we compile with xelatex + xeCJK. English body text is unaffected.
\usepackage[hidelinks]{hyperref}
\usepackage{etoolbox}
% Start every top-level section on a new page
\pretocmd{\section}{\clearpage}{}{}

\newcommand{\code}[1]{\texttt{#1}}
\newcommand{\kt}{\kappa_t}
\newcommand{\Nbar}{\bar N}
\newcommand{\Nhat}{\hat N}
\graphicspath{{../results/figures/}}

\title{\textbf{Market Concentration and the Flattening of the Phillips Curve}\\[2pt]
\large A Theory-Motivated, Semi-Structural State-Space Assessment of the HSA Mechanism}
\author{Satoshi Ichikawa}
\date{Draft --- August 2026}

\begin{document}
% Every quantitative claim in the body text is one of these macros. Numbers are
% not typed into the prose by hand -- see scripts/12_build_cpi_ppi_report.py.
% Unprefixed = interpolated (PCHIP) design; \Annual... = mixed-frequency design.
\input{../results/tables/quarterly_interpolated/result_macros.tex}
\input{../results/tables/quarterly_interpolated/bias_macros.tex}
\input{../results/tables/annual_q4/result_macros.tex}
\input{../results/tables/annual_q4/bias_macros.tex}
\input{../results/tables/shared/cross_design_macros.tex}
\input{../results/tables/shared/fit_comparison_macros.tex}
\input{../results/tables/shared/conditional_ml_macros.tex}
% Written by scripts/prior_decomposition_rho_delta.py, which is a separate
% diagnostic run-set. The file is \input FIRST and the ?? fallbacks come after:
% every macro file here uses \providecommand, which does not overwrite an
% existing definition, so a fallback placed first would silently win and the
% real numbers would never appear.
\IfFileExists{../results/tables/shared/prior_decomposition_macros.tex}{%
  \input{../results/tables/shared/prior_decomposition_macros.tex}}{}
\providecommand{\PriorDecompHeadlineDeltaPriorShare}{??}
\providecommand{\PriorDecompHeadlineArTwoPriorShare}{??}
\providecommand{\PriorDecompHeadlineInteractionShare}{??}
\providecommand{\PriorDecompCoreDeltaPriorShare}{??}
\providecommand{\PriorDecompCoreArTwoPriorShare}{??}
\providecommand{\PriorDecompCoreInteractionShare}{??}

\tableofcontents
\clearpage

\section{Why we estimate this equation: the HSA Phillips curve}
\label{sec:why}
Our estimating equation is the theory-motivated counterpart of the New Keynesian
Phillips curve derived by Fujiwara--Matsuyama (2026, \emph{Journal of Monetary
Economics}) under the HSA demand system. This section states their result, maps it to
what we take to data, and makes explicit how far the empirical specification departs
from the theory.

\subsection{The HSA Phillips curve}
A final-good producer aggregates a continuum of varieties through an HSA demand system
with market-share function $s(z)$, where $z=p/A$ is a firm's price normalized by the
single price aggregator $A$ (the reference price faced by firms, an inverse index of
competitive pressure). Intermediate firms set prices under a Rotemberg adjustment cost
$\chi$ and enter endogenously with a one-period time-to-build lag. Solving the firm's
problem, PPI inflation follows (their Eq.~21)
\begin{equation}
\hat\pi_t=\beta(1-\delta)\,\mathbb E_t\hat\pi_{t+1}
+\underbrace{\frac{\zeta(z)-1}{\chi}}_{\text{slope on marginal cost}}
\big(\hat W_t-\hat Z_{P,t}-\hat p_t\big)
-\underbrace{\frac{1}{\chi}\,\frac{1-\rho(z)}{\rho(z)}}_{\text{entry / cost-push}}\hat N_t,
\label{eq:fm}
\end{equation}
where $\hat N_t$ is the inverse HHI (number of firms), $\hat W_t-\hat Z_{P,t}-\hat p_t$ is the
inverse markup (the model's measure of real marginal cost), $\zeta(z)=1-s'(z)z/s(z)>1$
is the price-elasticity function (so the flexible-price markup is
$\mu=\zeta/(\zeta-1)$), and $\rho(z)$ is the pass-through rate.

\subsection{Two ways competition enters}
\paragraph{(1) The slope is a function of the firm count (structural flattening).}
The coefficient on marginal cost, $(\zeta(z)-1)/\chi$, depends on the price elasticity
$\zeta(z)$, which through $z=p/A$ and the symmetric-equilibrium condition $s(z)=1/N$ is
a function of the firm count $N$. Rising concentration---\emph{fewer firms}, higher
markups, lower pass-through $\rho(z)$---lowers this slope, so the Phillips curve
flattens \emph{structurally}. Linearizing the slope in the trend level N, 
$\bar N_t$ of the firm count gives
\begin{equation}
\kt=\kappa_0+\delta\,\bar N_t,
\end{equation}
so a \emph{positive} $\delta$ is the theory's flattening prediction: the slope moves
with the firm count and falls as concentration rises and $N$ declines.

\paragraph{(2) Endogenous entry is a cost-push term (estimation bias).} The second
term $-\tfrac{1}{\chi}\tfrac{1-\rho(z)}{\rho(z)}\hat N_t$ appears only away from CES.
It means that the cyclical component of firm entry acts as a cost-push shock through
strategic complementarity. Omitting it biases an empirical Phillips-curve estimate
(omitted-variable problem). We retain the cyclical term explicitly with its own
coefficient $\theta_t=\theta_0+\gamma\,\bar N_t$.

\paragraph{CES is the irrelevant case.} Under CES, $s(z)=\gamma_{\mathrm{CES}}\,z^{1-\theta}$,
so $\zeta(z)=\theta$ is constant and pass-through is $\rho(z)=1$. The entry term
vanishes and the slope is the constant $(\theta-1)/\chi$: market concentration neither
flattens the curve nor moves inflation. This is exactly our CES benchmark, and
departing from CES (the HSA models) is what can make $\delta$ and the $\hat N_t$ term
nonzero.

\subsection{From theory to the estimating equation}
Writing the activity/inverse-markup gap as $x_t$, adding a backward-looking inertia
term ($\alpha$) and a reduced-form shock $e_t$, and decomposing the observed firm
count into trend and cycle,
\begin{equation}
N^{\mathrm{obs}}_t=\bar N_t+\hat N_t+\nu_t,
\end{equation}
where $\bar N_t$ is a random-walk-with-drift trend, $\hat N_t$ is a stationary AR(2)
gap (truncated to the stationary region), and $\nu_t\sim\mathcal N(0,\sigma_N^2)$ is
measurement error, the HSA Phillips curve becomes the equation we estimate:
\begin{equation}
\boxed{\;\pi_t=\alpha\,\pi_{t-1}+(1-\alpha)\,\mathbb E_t\pi_{t+1}
+\kt\,x_t-\theta_t\,\hat N_t+e_t,\qquad
\kt=\kappa_0+\delta\,\bar N_t,\;\;\theta_t=\theta_0+\gamma\,\bar N_t.\;}
\label{eq:est}
\end{equation}
The trend $\bar N_t$ carries the structural-flattening channel (through $\delta$); the
cycle $\hat N_t$ carries the entry cost-push channel (through $\theta_t$).

\subsection{Distance from the theory}
\label{sec:distance}
We call this exercise \emph{semi-structural} rather than structural because the
specification departs from \eqref{eq:fm} in several ways. To avoid over-reading the
coefficients, we state each:
\begin{itemize}
\item the theory is about \emph{PPI} inflation, but there is no PPI expectation
series, so we pair observed PPI with a common external expectation series and the
expectation term is not a matched structural object;
\item the theoretical regressor is real marginal cost / the inverse markup, but the
main results use the negative unemployment gap.
\item we impose no cross-coefficient restrictions from the theory and approximate
$\kt$ as a \emph{linear} function of $\bar N_t$;
\item we add backward-looking inertia ($\alpha$), which is absent from \eqref{eq:fm};
\item expected inflation is an external survey measure (Section~\ref{sec:data});
\end{itemize}

\section{Data}
\label{sec:data}
\paragraph{Sample.} Quarterly, 1982Q1--2012Q4 ($T=124$) in the US.

\paragraph{Inflation.} We compare headline CPI, core CPI, and PPI. Expected inflation $\mathbb{E}\pi$ is
the same external survey series in all specifications. Because the Survey of
Professional Forecasters does not publish a PPI forecast, the PPI specifications use
the same expectation series as CPI. The PPI estimates are therefore not a structural
estimate matched through the expectation term, but a semi-structural external-validity
check that changes only the observed inflation index. A measurement mismatch with the
expectation series also remains for core CPI.

\paragraph{Activity gaps (five).} Negative unemployment gap, BN output gap, HP output gap
($100\times$ log real output HP-filtered, in the same $100$-log-point units as BN),
HP labor-share gap, and inverse markup. The unemployment measure enters as the \emph{negative} unemployment gap $x_t=-(u_t-u^*_t)$---labor-market tightness, positive in expansions and negative in downturns (e.g.\ $x_t\approx-4$ in 2009), co-moving positively with the output gaps. A positive slope $\kappa$ on $x_t$ is therefore a conventional downward-sloping Phillips curve.

\paragraph{Competition proxy, and why the main results are mixed-frequency.} The
competition series is Gustavo's firm-count measure, i.e.\ the inverse of the \emph{HHI of
U.S.\ listed firms}. It is \textbf{annual}: over 1982--2012 there are 31 observations, not
124.

The main results therefore use a \emph{mixed-frequency} observation scheme. The annual
value enters the state-space model in that year's Q4 and the intervening 93 quarters are
treated as \emph{missing}, so the sub-annual firm count is inferred rather than assumed
(Section~\ref{sec:mixed-freq}). Alongside it we report the conventional alternative, in
which the annual series is PCHIP-interpolated to a quarterly value and that interpolated
value is then used \emph{as data} in all 124 quarters (Section~\ref{sec:pchip}).

Earlier revisions of this paper took the interpolated version as the main result. We
reversed that after finding that the interpolation is not innocuous. It manufactures
within-year variation that does not exist, drives the estimated firm-count measurement
error down to $\sigma_N\approx\MainSigmaN$ (versus $\AnnualMainSigmaN$ under mixed
frequency), and through
that pins $\bar N_t+\hat N_t$ so tightly at all 124 quarters that only the \emph{split}
between trend and cycle is left free. The consequence is a near-exact posterior ridge
between the two states---posterior correlation $\MainStateCorr$ under interpolation against
$\AnnualMainStateCorr$ under mixed frequency---which is the identification pathology
described in
Section~\ref{sec:scope}. The interpolated design also forces the AR(2) cycle to a
near-unit root and hence to a $\MainArTwoPeriod$-quarter ``cycle'' that is a second trend
rather than a
cycle. None of this is visible if the interpolated series is simply taken as data.

The reassuring part is that $\delta$ itself barely moves between the two designs, so the
substantive result does not hinge on the choice; what changes is how much of the reported
precision is real.

\paragraph{Unit convention.} The default transform is
$(100\log N-\overline{100\log N})/10$, so one unit of $\bar N$ or $\hat N$ is a
ten-log-point change around the sample mean. The reported
$\delta,\theta,\theta_0,\gamma$ are all on this ten-log-point scale.

\begin{figure}[H]
\centering
\includegraphics[height=0.85\textheight,keepaspectratio]{shared/data_series.png}
\caption{Data. (a) Headline, core, and PPI inflation. (b) Real-activity gaps
(negative unemployment gap, BN and HP output gaps). (c) Competition measure
(estimation transform of the listed-firm count).}
\label{fig:data}
\end{figure}

\section{Econometric model and estimation}
\label{sec:model}

\subsection{Model variants}
The five models differ only in the restriction on the slope (Table~\ref{tab:models}).

\begin{table}[H]
\centering
\caption{Model variants.}
\label{tab:models}
\input{../results/tables/shared/model_variants.tex}
\end{table}

\paragraph{State blocks.} The bilinear term $-\gamma\bar N_t\hat N_t$ makes the joint
firm-count state non-linear-Gaussian, so \textbf{HSA full is estimated by a Particle
Gibbs sampler}: one conditional sequential-Monte-Carlo sweep updates the whole state
path (512 particles). This applies to \emph{every} HSA-full cell, under both the PCHIP
and the annual-Q4 observation scheme, so the PCHIP-versus-annual-Q4 comparison of
Section~\ref{sec:pchip} differs only in the observation scheme and not in the sampler. The table builder refuses to emit a table in which any HSA-full cell was
drawn by a different sampler.

Every other model draws the whole state $s_t=(\hat N_t,\hat N_{t-1},\bar N_t)'$ jointly
by exact Kalman/FFBS. That includes \textbf{HSA const-theta}: with $\gamma=0$ the
inflation observation is linear in the joint state, so the model is linear-Gaussian and
admits an exact joint draw. Earlier revisions sampled it with the same two alternating
blocks as HSA full; because $\bar N_t$ and $\hat N_t$ are almost perfectly negatively
correlated in the posterior under interpolation ($\MainStateCorr$: the firm-count equation
pins their sum but
barely their split), that blocking moved their common level with an integrated
autocorrelation time of order $10^3$ sweeps and did not mix. This is a purely
computational change---it removes a sampler pathology and creates no identification that
the data do not contain. Coefficients, variances, priors, chains, iterations, burn-in and
thinning are unchanged throughout. 

\subsection{The inflation equation actually estimated}
To put the theoretical equation \eqref{eq:est} directly in regression form, we move
expected inflation to the left-hand side and define
\begin{equation}
y_t\equiv\pi_t-E_t\pi_{t+1},\qquad
a_t\equiv\pi_{t-1}-E_t\pi_{t+1}
\label{eq:ya-def}
\end{equation}
Making the unforecastable component of the activity variable explicit, the estimating
system is
\begin{align}
y_t
 &=\alpha a_t+(\kappa_0+\delta\bar N_t)x_t
   -(\theta_0+\gamma\bar N_t)\hat N_t
   +\lambda_{e\zeta}\zeta_t+\eta_t, \label{eq:estimated-inflation}\\
x_t&=\phi_1x_{t-1}+\zeta_t,\qquad
\eta_t\perp\zeta_t. \label{eq:activity-ar1}
\end{align}
That is, the original NKPC shock decomposes as $e_t=\lambda_{e\zeta}\zeta_t+\eta_t$.
Estimating $\lambda_{e\zeta}$ allows a correlation between the contemporaneous activity
shock and the inflation shock, rather than forcing a simple regression that treats
$x_t$ as exogenous. The saved correlation coefficient is
\begin{equation}
\rho_{e\zeta}
=\frac{\lambda_{e\zeta}\sigma_\zeta}
{\sqrt{\lambda_{e\zeta}^2\sigma_\zeta^2+\sigma_\eta^2}}
\end{equation}
The main HSA steady sets $\theta_0=\gamma=0$; HSA dynamic sets
$\delta=\gamma=0$.

For numerical stability the implementation updates internal coefficients on the
rescaled regressors $x_t/100$ and $x_t\bar N_t/100$ for $\kappa_0$ and $\delta$; the
coefficients are rescaled back by $100$ on saving, so the $\kappa_0,\delta,\kt$ in the
text, tables, and figures are all in the physical units of
\eqref{eq:estimated-inflation}. They must not be multiplied by $100$ or divided by
$100$ again.

\subsection{State-space representation splitting the firm count into trend and cycle}
The observation and transition equations for the firm count are
\begin{align}
N_t^{\mathrm{obs}}&=\bar N_t+\hat N_t+\nu_t,
&\nu_t&\sim N(0,\sigma_N^2), \label{eq:N-measurement}\\
\hat N_t&=\rho_1\hat N_{t-1}+\rho_2\hat N_{t-2}+u_t,
&u_t&\sim N(0,\sigma_u^2), \label{eq:N-gap}\\
\bar N_t&=n+\bar N_{t-1}+\varepsilon_t,
&\varepsilon_t&\sim N(0,\sigma_\varepsilon^2) \label{eq:N-trend}
\end{align}
where $\bar N_t$ is the low-frequency trend with drift $n$ and $\hat N_t$ is a
stationary AR(2) cycle. Rather than a two-step procedure that first filters the
observed firm count into components and then substitutes them into the NKPC, we extract
both states using the firm-count equation and the inflation equation within the same
Gibbs iteration.

\paragraph{Two observation frequencies; the main results use the first.}
\label{sec:mixed-freq}
The firm count is annual, so the \emph{mixed-frequency} design places the annual value in
each year's Q4 and treats Q1--Q3 as missing,
\begin{equation}
N_t^{\mathrm{obs}}=
\begin{cases}
(100\log N_{y(t)}^{\mathrm{annual}}-\bar c)/10,&t\in\mathcal T_N\;(\mathrm{Q4}),\\
\mathrm{missing},&t\notin\mathcal T_N\;(\mathrm{Q1--Q3}),
\end{cases}
\qquad |\mathcal T_N|=31,
\label{eq:annual-q4-observation}
\end{equation}
where $\bar c$ is a fixed centering constant, shared with the interpolated series so the
coefficient units are comparable across designs. The alternative \emph{interpolated}
design instead PCHIP-interpolates the annual series to quarterly and applies
\eqref{eq:N-measurement} at all 124 quarters; Section~\ref{sec:pchip} reports it and
documents what the interpolation costs. In Q4 the Kalman update uses two observation
rows---the firm-count equation and the inflation equation---and in Q1--Q3 only the one
inflation row. Because the state transition and prediction continue at every quarter,
the Q1--Q3 $\bar N_t,\hat N_t$ are not simple interpolations but FFBS draws conditioned
on the annual anchors, the state transition, and the inflation equation, updated every
iteration. The inverse-gamma update of $\sigma_N^2$ likewise uses only the 31 finite
Q4 residuals.

For the main HSA steady the state is
\begin{equation}
s_t=(\hat N_t,\hat N_{t-1},\bar N_t)',\qquad
s_t=c+Fs_{t-1}+w_t
\end{equation}
with
\begin{equation}
c=\begin{bmatrix}0\\0\\n\end{bmatrix},\quad
F=\begin{bmatrix}\rho_1&\rho_2&0\\1&0&0\\0&0&1\end{bmatrix},\quad
Q=\operatorname{Var}(w_t)
=\begin{bmatrix}\sigma_u^2&0&0\\0&0&0\\0&0&\sigma_\varepsilon^2\end{bmatrix}.
\label{eq:state-matrices}
\end{equation}
Stacking the firm-count observation with the inflation equation net of its known part,
\begin{equation}
y_t^\pi=y_t-\alpha a_t-\kappa_0x_t
-\lambda_{e\zeta}\zeta_t
\end{equation}
gives
\begin{equation}
\underbrace{\begin{bmatrix}N_t^{\mathrm{obs}}\\y_t^\pi\end{bmatrix}}_{z_t}
=
\underbrace{\begin{bmatrix}1&0&1\\0&0&\delta x_t\end{bmatrix}}_{H_t}s_t
+v_t,\qquad
\operatorname{Var}(v_t)=
\begin{bmatrix}\sigma_N^2&0\\0&\sigma_\eta^2\end{bmatrix}.
\label{eq:steady-measurement-matrix}
\end{equation}
Conditional on the coefficients and variances this is a linear-Gaussian state-space
model, so $s_{0:T}$ can be drawn jointly by an exact Kalman filter /
forward-filtering backward-sampling (FFBS).

The initial state starts from the fixed prior
\begin{equation}
s_0=(\hat N_0,\hat N_{-1},\bar N_0)'\sim N(m_0,P_0),\qquad
m_0=(0,0,0)',\quad P_0=10I_3
\label{eq:initial-state}
\end{equation}
We do not substitute $N_0^{\mathrm{obs}}$ into $m_0$; the period-0 firm count and
inflation enter the ordinary Kalman observation update. The first AR(2) transition also
actually uses the prior-updated $\hat N_{-1}$. This avoids double-using the initial
observation or dropping the first AR(2) likelihood.

\subsection{Conditional posteriors of the coefficients and variances}
Take the main model as an example, with $\beta=(\alpha,\kappa_0,\delta)'$,
$X_t=(a_t,x_t,x_t\bar N_t)$, and $y_t^*=y_t-\lambda_{e\zeta}\zeta_t$. Under the
independent normal prior $\beta\sim N(b_0,V_0)$, we update jointly from
\begin{align}
V_1&=\left(V_0^{-1}+\frac{X'X}{\sigma_\eta^2}\right)^{-1},\\
b_1&=V_1\left(V_0^{-1}b_0+\frac{X'y^*}{\sigma_\eta^2}\right),\\
\beta\mid\text{rest}&\sim N(b_1,V_1) \label{eq:beta-posterior}
\end{align}
The implementation uses the $100$-rescaled $X$ and coefficients noted above, but the
formula is identical. The baseline estimation imposes no sign restriction, so both
$\delta$ and the end-of-sample $\kt$ can be negative.

We update $\lambda_{e\zeta}$ from the normal conditional posterior of the regression of
$e_t$ on $\zeta_t$. We update $\phi_1$ using not only \eqref{eq:activity-ar1} but also
the information in $e_t=\lambda_{e\zeta}(x_t-\phi_1x_{t-1})+\eta_t$. We draw
$\rho_1,\rho_2$ from a normal conditional posterior truncated to the AR(2) stationary
region,
\begin{equation}
|\rho_2|<1,\qquad \rho_1+\rho_2<1,\qquad \rho_2-\rho_1<1
\end{equation}
The fallback count and rejection rate when no proposal is accepted within 2{,}000 draws
are saved to each run's metadata.

The variance priors are $\sigma^2\sim\operatorname{IG}(a,b)$ (density kernel
$(\sigma^2)^{-(a+1)}\exp(-b/\sigma^2)$), updated from the inverse-gamma conditional
posterior that adds the residual sum of squares. The concrete baseline priors are in
Table~\ref{tab:baseline-priors}.

\begin{table}[H]
\centering
\caption{Baseline priors for the main HSA steady model. Normal distributions show mean
and standard deviation; inverse-gamma distributions show $(a,b)$. The coefficient
priors are common to the other models.}
\label{tab:baseline-priors}
\small
\input{../results/tables/shared/baseline_priors.tex}
\end{table}

\noindent For HSA dynamic the shock covariance takes
$\Sigma\sim\operatorname{IW}(8,S_0)$ with $S_0=\operatorname{diag}(3,3,0.06,0.03)$.
Under the default $(e,\zeta)$-only, the corresponding marginals of this prior apply
to the $(e,\zeta)$ block and the two state variances.

\subsection{What the Kalman filter and FFBS do}
Write the predicted mean and variance at time $t$ as $a_t^s,P_{t|t-1}$ and the filtered
mean and variance as $m_t,P_t$. The forward update for
\eqref{eq:state-matrices}--\eqref{eq:steady-measurement-matrix} is
\begin{align}
a_t^s&=c+Fm_{t-1},&P_{t|t-1}&=FP_{t-1}F'+Q,\\
q_t&=z_t-H_ta_t^s,&S_t&=H_tP_{t|t-1}H_t'+R,\\
K_t&=P_{t|t-1}H_t'S_t^{-1},&m_t&=a_t^s+K_tq_t
\end{align}
The covariance update uses the numerically stable Joseph form. After drawing the
endpoint $s_T\sim N(m_T,P_T)$, we draw the smoothed state path jointly with
\begin{align}
J_t&=P_tF'P_{t+1|t}^{-1},\\
s_t\mid s_{t+1},z_{1:T}
&\sim N\!\left(m_t+J_t(s_{t+1}-c-Fm_t),
P_t-J_tP_{t+1|t}J_t'\right)
\end{align}
for $t=T-1,\ldots,0$. Thus the plotted $\bar N_t,\hat N_t,\kt$ carry the posterior
uncertainty not only of the coefficients but also of the state decomposition.

\subsection{One Gibbs iteration and the differences across models}
One iteration of HSA steady is, in order:
\begin{enumerate}
\item update $(\alpha,\kappa_0,\delta)$ from the current $\bar N_t$ via
\eqref{eq:beta-posterior};
\item update $\lambda_{e\zeta}$ and $\phi_1$, and recompute $\zeta_t,\eta_t$;
\item update $\sigma_\zeta^2,\sigma_\eta^2$;
\item update $(\rho_1,\rho_2)$ and $\sigma_u^2$;
\item update $n$ and $\sigma_\varepsilon^2$, then $\sigma_N^2$;
\item update the entire $s_{0:T}$ by FFBS using the updated coefficients and variances.
\end{enumerate}
CES is a conjugate regression benchmark with no latent firm-count state.\par
HSA dynamic handles the covariance of $r_t=(e_t,\zeta_t,u_t,\varepsilon_t)'$;
under the default $(e,\zeta)$-only it updates the $2\times2$ block of
$(e_t,\zeta_t)$ and $\sigma_u^2,\sigma_\varepsilon^2$ directly on the constrained
parameter space, rather than zeroing the off-diagonal of an unconstrained
inverse-Wishart draw after the fact.

HSA full contains the state-by-state product $-\gamma\bar N_t\hat N_t$, so the joint
firm-count state $(\bar N_t,\hat N_t)$ is not linear-Gaussian at once. HSA full is therefore updated by the joint \emph{Particle Gibbs} step detailed in
Section~\ref{sec:pg}---a conditional sequential-Monte-Carlo update of the whole state
path (512 particles)---which replaces the two alternating exact-FFBS blocks used
previously. HSA const-theta ($\gamma=0$) is a different case: without the bilinear term
the inflation observation loads linearly on the joint state,
$\big[-\theta_0,\,0,\,\delta x_t\big]s_t$, so the model is linear-Gaussian and the whole
path is drawn in a single exact FFBS sweep---the same routine HSA steady uses, of which
HSA steady is the $\theta_0=0$ case. In both models the weak signal of the cyclical state
$\hat N_t$, which must be identified as a separate coefficient, keeps identification
weaker than for the other models; neither the joint Particle Gibbs update nor the exact
joint FFBS removes that weak identification, and neither is claimed to.

\subsection{The Particle Gibbs state update (HSA full)}
\label{sec:pg}
For HSA full the bilinear term $-\gamma\bar N_t\hat N_t$ in
\eqref{eq:estimated-inflation} makes the joint state $s_t=(\hat N_t,\hat N_{t-1},\bar
N_t)'$ non-linear-Gaussian, so the exact FFBS blocks used above do not apply jointly.
We instead draw the whole path $s_{0:T}$ with one \emph{Particle Gibbs} sweep---a
conditional sequential-Monte-Carlo step that leaves the exact conditional
$p(s_{0:T}\mid\theta,y,N^{\mathrm{obs}})$ invariant for any number of particles $P$
(Andrieu, Doucet and Holenstein, 2010).

\paragraph{Proposal and weights.} Particles are propagated from the state transition
\eqref{eq:N-gap}--\eqref{eq:N-trend} (a bootstrap proposal): given the ancestor
$s_{t-1}^{(i)}$,
\begin{equation}
\hat N_t^{(i)}=\rho_1\hat N_{t-1}^{(i)}+\rho_2\hat N_{t-2}^{(i)}+u_t^{(i)},\qquad
\bar N_t^{(i)}=n+\bar N_{t-1}^{(i)}+\varepsilon_t^{(i)},\qquad
u_t^{(i)}\sim N(0,\sigma_u^2),\;\varepsilon_t^{(i)}\sim N(0,\sigma_\varepsilon^2).
\end{equation}
Each particle is weighted by \emph{both} observation equations---the inflation equation
\eqref{eq:estimated-inflation} and the firm-count measurement \eqref{eq:N-measurement}:
\begin{align}
\log w_t^{(i)}&=
-\frac{\big(y_t-\mu_t^{(i)}\big)^2}{2\sigma_\eta^2}
-\frac{\big(N_t^{\mathrm{obs}}-\bar N_t^{(i)}-\hat N_t^{(i)}\big)^2}{2\sigma_N^2},
\label{eq:pg-weight}\\
\mu_t^{(i)}&=\alpha a_t+\big(\kappa_0+\delta\bar N_t^{(i)}\big)x_t
-\big(\theta_0+\gamma\bar N_t^{(i)}\big)\hat N_t^{(i)}+\lambda_{e\zeta}\zeta_t,
\end{align}
with $y_t=\pi_t-\mathbb E_t\pi_{t+1}$, $a_t=\pi_{t-1}-\mathbb E_t\pi_{t+1}$ and
$\zeta_t=x_t-\phi_1 x_{t-1}$ (constants common to all particles at $t$ cancel under
normalisation). Normalised weights $W_t^{(i)}=w_t^{(i)}/\sum_j w_t^{(j)}$ are formed with
a log-sum-exp for stability, and $\mathrm{ESS}_t=1/\sum_i (W_t^{(i)})^2$.

\paragraph{Conditional SMC.} Let $s_{0:T}^{\mathrm{ref}}$ be the state path from the
previous MCMC iteration, pinned to particle slot $1$.
\begin{enumerate}
\item $t=0$: set $s_0^{(1)}=s_0^{\mathrm{ref}}$ and draw $s_0^{(i)}\sim N(m_0,P_0)$ for
$i\ge2$ from the initial prior \eqref{eq:initial-state}; weight by \eqref{eq:pg-weight}.
\item $t=1,\dots,T$: resample ancestors $a_t^{(i)}\sim\mathrm{Categorical}(W_{t-1})$ for
$i\ge2$ and set $a_t^{(1)}=1$, so the reference keeps its lineage; propagate
$s_t^{(i)}$ from the transition above given $s_{t-1}^{(a_t^{(i)})}$, with
$s_t^{(1)}=s_t^{\mathrm{ref}}$; weight by \eqref{eq:pg-weight}. The pair
$(\hat N_t,\hat N_{t-1})$ in $s_t$ preserves the AR(2) genealogy.
\item Draw a terminal index $J\sim\mathrm{Categorical}(W_T)$ and trace its ancestors
backward to return one joint path $s_{0:T}$, which is the new state draw and the
reference for the next iteration.
\end{enumerate}

\subsection{Iterations, saved draws, and convergence criteria}
Each cell uses two independent chains, with a \emph{total} of 12{,}000 iterations per
chain. We discard the first 4{,}000 as burn-in and save every 5th of the remaining
8{,}000, giving 1{,}600 draws per chain and 3{,}200 draws per cell for diagnostics and
summaries. The main analysis is 75 cells: the baseline 45 (3 price indices $\times$ 3
activity measures $\times$ 5 models) plus the 30 weak/tight-prior cells. Separately, we
estimate 2 inverse-markup CES/HSA-dynamic cells with the same settings.

\paragraph{What ``converged'' means here.} The threshold is maximum $\hat R\leq1.01$ and
minimum bulk ESS $\geq400$, but a single run has several distinct groups of quantities and
they do not pass or fail together. We therefore compute and report the rule on three
groups separately:
\begin{enumerate}
\item \textbf{scalar parameters} --- every stored scalar: $\alpha$, the slope terms, the
entry terms, $\rho_1,\rho_2$, the trend drift $n$, $\phi_1$, $\lambda_{e\zeta}$ and
\emph{every} estimated variance;
\item \textbf{latent-state paths} --- $\bar N_t$ and $\hat N_t$, maximum $\hat R$ and
minimum ESS along the path;
\item \textbf{derived paths} --- $\kt$ and, where estimated, $\theta_t$.
\end{enumerate}
The $\dagger$ mark in the coefficient tables is the \emph{scalar} rule. A cell marked
``OK (coef)'' passes on the coefficients but not on the latent-state paths; ``OK'' passes
on all three. We state this explicitly because the two are routinely conflated: marginal
coefficient chains can mix well while the joint posterior including the states has not,
and reporting only the former would overstate what has been established.
Table~\ref{tab:group-convergence} gives the group-by-group numbers for every model.

Under the main mixed-frequency design, \AnnualReportWarningCount{} of the
\AnnualReportRunCount{} runs fail the scalar rule; under the interpolated design the
figures are \ReportWarningCount{} of \ReportRunCount{}. The two designs fail for different
reasons, and the difference is informative rather than incidental. Under interpolation the
worst-mixing scalar is the trend drift $n$ in \DriftWorstCount{} of the \HsaFailCount{}
failing HSA cells---the level ridge
of Section~\ref{sec:scope} surfacing in the parameter block. Under mixed frequency it is
the AR(2) block $(\rho_1,\rho_2)$ in \AnnualArTwoWorstCount{} of \AnnualHsaCellCount{}
cells, which for HSA steady is a nuisance
object because $\theta=0$ keeps $\hat N_t$ out of the inflation equation. The parameters
the conclusions rest on fail far less often than the run as a whole, under either design:
$\delta$ in \AnnualDeltaFailCount{} of \AnnualHsaCellCount{} cells and $\kappa_0$ in
\AnnualKappaZeroFailCount{} under mixed frequency, against \DeltaFailCount{} and
\KappaZeroFailCount{} under interpolation. The $95\%$ credible intervals are the $2.5\%$ and $97.5\%$
quantiles of the saved draws; the plotted points are posterior means.

\section{Main results}
\label{sec:main}
All results in this section use the \textbf{mixed-frequency} firm-count design: the 31
annual observations enter in their own Q4 and the 93 intervening quarters are inferred
inside the state-space model (Section~\ref{sec:mixed-freq}). Table~\ref{tab:headline}
gives the main-model (HSA steady) estimates for every price index and activity measure,
so each result can be read against the specification it comes from. The main
specification---\textbf{core CPI, negative unemployment gap}---is in bold; a positive
$\delta$ means the slope flattens as the firm count falls.
\input{../results/tables/annual_q4/headline_results.tex}

\paragraph{Reading the convergence column.} Every cell in Table~\ref{tab:headline}
carries a $\dagger$, and it would be a mistake to read that as ``$\delta$ is not
estimated reliably here.'' The blanket flag is set by whichever scalar mixes worst, and
under mixed frequency that is the AR(2) block $(\rho_1,\rho_2)$ in \AnnualArTwoWorstCount{}
of the \AnnualHsaCellCount{} HSA cells.
For \textbf{HSA steady} the AR(2) block governs $\hat N_t$, which---because $\theta=0$ in
this model---does not enter the inflation equation at all. The coefficient the
conclusions rest on is $\delta$, and its own diagnostics are in the penultimate column:
$\delta$ alone clears $\hat R\leq1.01$ and bulk ESS $\geq400$ in
\AnnualSteadyDeltaPassCount{} of the
\AnnualSteadyDeltaCellCount{} HSA-steady baseline cells, at a median bulk ESS of
\AnnualSteadyDeltaMedianEss{}. That is a statement about one coefficient, not about the
run: no HSA cell in this design passes the whole-run rule, which is applied to all
\ScalarParameterCount{} scalars at once.
Section~\ref{sec:scope} states which failures do and do not bear on the conclusions.

Table~\ref{tab:model-comp} then places the main model against the other model variants
(CES, HSA dynamic, HSA const-theta, HSA full) for the negative unemployment gap. CES and
HSA dynamic hold the slope constant, so they have no $\delta$; the HSA models with a
firm-count-dependent slope estimate $\delta$, and HSA steady is the main model. Full
coefficients ($\theta$, $\gamma$, and the $\kappa_t$ path) for every activity measure are
in Appendix~\ref{app:model-tables}.
\input{../results/tables/annual_q4/model_comparison_unemp.tex}

\subsection{The negative-unemployment-gap slope flattens as the firm count falls}
For the baseline HSA steady on the main sample, the core-CPI negative-unemployment-gap
specification gives
\[
\delta=\AnnualCoreUnempDelta,
\]
and the Savage--Dickey Bayes factor against $\delta=0$ is
$\mathrm{BF}_{10}=\AnnualCoreUnempDeltaBF$. Because $\bar N_t$ falls over the sample, the
implied slope $\kt$ declines monotonically from $\AnnualCoreUnempKappaStart$ to
$\AnnualCoreUnempKappaEnd$: as the firm count falls (concentration rises), the Phillips
curve flattens, the sign the HSA mechanism predicts. Headline CPI also gives
$\delta=\AnnualHeadlineUnempDelta$
($\mathrm{BF}_{10}=\AnnualHeadlineUnempDeltaBF$), so the CPI result for the negative
unemployment gap is robust to the inflation index. Figure~\ref{fig:kt_unemp} plots the
implied $\kt$ path; its end-of-sample band includes zero, so the accurate statement is a
flattening to near zero rather than a strictly positive residual slope.

\begin{figure}[H]
\centering
\includegraphics[width=0.62\linewidth]{annual_q4/kappa_t_unemployment_by_inflation.png}
\caption{Time-varying slope $\kt$ for the negative unemployment gap under the
mixed-frequency design (posterior mean and $95\%$ band). The band includes the
uncertainty from inferring the 93 unobserved quarterly firm counts. Under CPI it declines
to near zero by the end of the sample, whereas the PPI band is wide.}
\label{fig:kt_unemp}
\end{figure}

The estimated Phillips-curve slope in the core-CPI specification declines from
$\AnnualCoreUnempKappaStart$ at the beginning of the sample to approximately zero by the
end. This change is
economically meaningful. For example, a four-point deterioration in the negative
unemployment gap ($x_t=-4$) implies an inflation response of about
$\AnnualCoreUnempFourPointResponse$ percentage
point at the 1982 slope, but almost no response at the 2012 slope.

\begin{figure}[H]
\centering
\includegraphics[width=0.88\linewidth]{annual_q4/competition_decomposition_hsa_steady_core_unemployment.png}
\caption{Posterior decomposition of the firm count $N_t=\bar N_t+\hat N_t$ under the main
HSA steady (negative unemployment gap, core CPI), mixed-frequency design. Lines are
posterior means, bands are $95\%$ credible intervals. The quarterly states between the Q4
annual observations are inferred, not interpolated, so the bands widen between anchors.}
\label{fig:main-state-decomposition}
\end{figure}

\section{The CES slope bias and the constant-slope benchmark}
\label{sec:ces-hsa-bias}\label{sec:bench}

\subsection{Theoretical prediction and its empirical counterpart}
Fujiwara--Matsuyama (2026) show that because endogenous entry generates a cost-push
term through strategic complementarity, a CES-type equation that regresses inflation on
real marginal cost alone under-estimates the NKPC slope. In our HSA-dynamic
representation, moving expected inflation to the left-hand side gives
\begin{equation}
y_t=\alpha a_t+\kappa_{\mathrm{HSA}}x_t+q_t+e_t,
\qquad q_t=-\theta\hat N_t,
\label{eq:bias-hsa}
\end{equation}
where $y_t=\pi_t-\mathbb E_t\pi_{t+1}$ and $a_t=\pi_{t-1}-\mathbb E_t\pi_{t+1}$. When
the CES equation omits $q_t$,
\begin{equation}
\operatorname{Bias}(\hat\kappa_{\mathrm{CES}})
=\operatorname{plim}(\hat\kappa_{\mathrm{CES}}-\hat\kappa_{\mathrm{HSA}})
=\frac{\tilde x'\tilde q}{\tilde x'\tilde x},
\label{eq:bias-ovb}
\end{equation}
where tildes denote residuals after removing $a_t$ and the activity shock
$\zeta_t=x_t-\phi_1x_{t-1}$. The theory predicts \eqref{eq:bias-ovb} is negative, i.e.\
$\kappa_{\mathrm{CES}}-\kappa_{\mathrm{HSA}}<0$.

To keep this correspondence sharp, separately from the slack measures of the main
results we use the BN-decomposed inverse markup as a proxy for real marginal cost and
re-estimate CES and HSA dynamic with the same sample, priors, and MCMC settings. For
each HSA posterior draw we form $q_t^{(s)}=-\theta^{(s)}\hat N_t^{(s)}$ and directly
compute the bias in \eqref{eq:bias-ovb}. Because the CES--HSA difference is a comparison
across separately estimated posteriors, we build a Monte Carlo distribution of the
difference by randomly permuting the HSA draws with a fixed seed.

\subsection{Estimation results}
For the inverse-markup specification that corresponds directly to the theory,
\begin{align*}
\kappa_{\mathrm{CES}}&=\AnnualBiasDirectCES,&
\kappa_{\mathrm{HSA}}&=\AnnualBiasDirectHSA,\\
\kappa_{\mathrm{CES}}-\kappa_{\mathrm{HSA}}&=\AnnualBiasDirectDifference,&
\Pr(\kappa_{\mathrm{CES}}-\kappa_{\mathrm{HSA}}<0)&=\AnnualBiasDirectProbability,\\
B_{\mathrm{HSA}}&=\AnnualBiasDirectOVB,&
\Pr(B_{\mathrm{HSA}}<0)&=\AnnualBiasDirectOVBProbability
\end{align*}
The CES--HSA difference does not show a negative bias. The omitted-variable bias
implied by HSA itself also has a point estimate of the sign opposite to the theory
(positive), with a credible interval that includes zero.

\begin{table}[H]
\centering
\caption{CES--HSA-dynamic slope comparison and the omitted-variable bias implied by
HSA. Each interval is a $95\%$ credible interval. $\dagger$ marks that at least one of
CES or HSA dynamic is outside the convergence criteria.}
\label{tab:ces-hsa-bias}
\tiny
\resizebox{\linewidth}{!}{\input{../results/tables/annual_q4/ces_hsa_kappa_bias.tex}}
\end{table}

\paragraph{Verdict.} The inverse-markup row passes the convergence criteria for both
models, but the slope intervals are wide and neither the CES--HSA difference nor the
within-model OVB supports the negative sign the theory predicts. The three auxiliary
rows using the negative unemployment gap give the same conclusion. This does not amount to a
refutation of the theory: the inverse markup is an aggregate series rather than
firm-level real marginal cost, and HSA dynamic is not a joint estimate of the full
theoretical model. What we can say from this paper is only that \emph{in the current
data and semi-structural model, we do not obtain an estimate corresponding to the
theory paper's prediction of a negative bias in the CES slope}.

As a benchmark for the time-varying result of Section~\ref{sec:main}, we also estimate
the constant-slope models. The full model-by-model coefficient tables for the
negative unemployment gap, and the corresponding prior-sensitivity table, are collected in
Appendix~\ref{app:model-tables}.
\paragraph{CES vs HSA (constant $\kappa$) } CES is the constant-slope benchmark. Among the
converged cells, the core-CPI $\kappa$ interval sits on the positive side, while
headline CPI and PPI include zero. CES measures the existence of the average curve but
does not explain the flattening. The prior-vs-posterior of the slope
(Figure~\ref{fig:kappa_pp}) shows CES's constant $\kappa$ and HSA steady's
slope at average competition $\kappa_0$, both sharply updated by the data and sitting
just above zero.

\begin{figure}[H]
\centering
\includegraphics[width=0.55\linewidth]{quarterly_interpolated/pp_kappa_ces_vs_steady.png}
\caption{Prior vs.\ posterior of the Phillips-curve slope for the negative unemployment gap:
CES's constant slope $\kappa$ and HSA steady's slope at average competition
$\kappa_0$, against their common prior $\mathcal N(0.1,0.2)$. Both are sharply updated
by the data and sit just above zero; the two models agree on the slope level, so what
distinguishes them is whether that slope varies with the firm count
(Section~\ref{sec:mc}).}
\label{fig:kappa_pp}
\end{figure}

\section{Does the time-varying slope improve fit over the constant slope?}
\label{sec:mc}
We compare CES (constant $\kappa$) against HSA steady (slope varying with the firm
count) by two criteria: predictive fit (Table~\ref{tab:fit-comparison}) and the
Savage--Dickey Bayes factor for $\delta\neq 0$ (Table~\ref{tab:mc}).

\paragraph{Fit, by four scores at once.} Earlier revisions reported only an
\emph{in-sample plug-in} score. It evaluates the inflation equation at posterior-mean
parameters and states (bars denote posterior means),
\begin{equation}
\hat\pi_t=\bar\alpha\,\pi_{t-1}+(1-\bar\alpha)\,\mathbb{E}_t\pi_{t+1}+\bar\kt\,x_t,
\qquad \bar\kt=\bar\kappa_0+\bar\delta\,\bar N_t
\;\;(\text{CES: }\bar\kt=\bar\kappa),
\end{equation}
and scores the residuals $r_t=\pi_t-\hat\pi_t$ by a Gaussian log-likelihood at the
posterior-mean shock variance,
\begin{equation}
S=-\frac{T}{2}\log\!\big(2\pi\bar\sigma_e^2\big)-\frac{1}{2\bar\sigma_e^2}\sum_{t=1}^{T}r_t^2 .
\end{equation}
This reuses the estimation sample and collapses the posterior to a point, so it is
neither an out-of-sample forecast nor a posterior predictive density.

We now report it \emph{alongside} three proper scores rather than in place of them:
the prequential one-step-ahead log predictive density
$\mathrm{LPD}_1=\sum_t\log p(\pi_t\mid\pi_{1:t-1},x,N)$, integrated over the posterior;
WAIC; and PSIS-LOO. Keeping the plug-in score visible is the point of the table---it lets
the reader see how much a crude criterion overstates the case:

\input{../results/tables/annual_q4/fit_comparison.tex}

For HSA steady against CES the plug-in score gives $\FitCorePluginDiff$ on core CPI where
$\mathrm{LPD}_1$ gives $\FitCoreLpdDiff$ and WAIC and LOO give $\FitCoreWaicDiff$ and
$\FitCoreLooDiff$; on headline CPI
it gives $\FitHeadlinePluginDiff$ against $\FitHeadlineLpdDiff$, $\FitHeadlineWaicDiff$
and $\FitHeadlineLooDiff$. The plug-in score overstates the
improvement by a factor of roughly $\FitPluginOverstateMin$ to $\FitPluginOverstateMax$
against the prequential score, which is what one expects from a
criterion that rewards a model for the extra latent states it fitted to this same sample.
The \emph{ordering} nevertheless survives: on both CPI indices the time-varying slope
improves every score, on PPI no score meaningfully separates the models (LPD$_1$
$\FitPPILpdDiff$, WAIC $\FitPPIWaicDiff$,
LOO $\FitPPILooDiff$), and HSA dynamic---the entry channel without a time-varying slope---never
beats HSA steady. Two caveats. First, the LOO column is unreliable wherever the maximum Pareto
$\hat k$ exceeds $0.7$, which is \FitLooFlaggedCount{} of the \FitLooCellCount{} cells
(the $\hat k$ column of Table~\ref{tab:fit-comparison} marks which).
That includes the CES baseline for the core-CPI comparison, so the $\FitCoreLooDiff$ LOO difference
quoted above rests on one number that PSIS itself flags; the WAIC difference of
$\FitCoreWaicDiff$,
which does not depend on importance sampling, is the more trustworthy of the two, and it
agrees. The PPI cells should simply not be read off the LOO column. Second, all four
scores are computed on the estimation sample---LPD$_1$ is prequential in the sense that
each $\pi_t$ is scored against a predictive density formed without $\pi_t$, but the
parameters were estimated on the full sample---so none of this is a genuine
out-of-sample exercise.

\paragraph{Bayes factor (Savage--Dickey), stated precisely.} We test $\delta=0$
\emph{within the HSA state model}: $M_1$ is HSA steady with $\delta$ free and
$M_0$ is the same model with $\delta=0$. The Savage--Dickey identity gives
\begin{equation}
\mathrm{BF}_{10}=\frac{p(\delta=0)}{\pi(\delta=0\mid y)},\qquad
p(\delta=0)=\frac{1}{\sigma_\delta\sqrt{2\pi}}=\frac{1}{0.02\sqrt{2\pi}}\approx 19.95
\;\;(\delta\sim\mathcal N(0,\sigma_\delta^2)),
\end{equation}
The posterior ordinate is computed \emph{without density estimation}. The Gibbs update
for the inflation coefficients is an exact Gaussian conditional, so the marginal
ordinate is its Rao--Blackwellised average over the retained draws of everything else,
$\psi=(\bar N_{0:T},\hat N_{0:T},\sigma_\eta^2,\lambda_{e\zeta},\phi_1)$:
\begin{equation}
\hat\pi(\delta=0\mid y)=\frac{1}{S}\sum_{s=1}^{S}
\mathcal N\!\big(0;\;b_1^{(s)}[\delta],\,V_1^{(s)}[\delta,\delta]\big),
\qquad \psi^{(s)}\sim p(\psi\mid y),
\label{eq:sddr-rb}
\end{equation}
with $b_1^{(s)},V_1^{(s)}$ the conditional moments of \eqref{eq:beta-posterior}
evaluated at $\psi^{(s)}$. This matters because $\delta=0$ lies about three posterior
standard deviations outside the sampled range---no retained draw of $\delta$ is
negative---so a kernel density evaluated there extrapolates its own tail rather than
measuring the posterior: a Gaussian kernel estimate at that point moves by orders of
magnitude as its bandwidth is scaled around Scott's
rule, whereas \eqref{eq:sddr-rb} draws on \AnnualSddrEffectiveTerms{} effectively
contributing terms out
of \AnnualSddrTotalTerms{} (the three largest carry $\AnnualSddrTopThreeShare\%$ of the
sum) and has a Monte Carlo standard
error of about $\AnnualSddrMcsePercent\%$.

Three cautions remain. (i) This is \emph{not} a CES-vs-HSA marginal-likelihood ratio:
CES and HSA also differ in the latent-$N$ block, which the SDDR holds fixed, so we keep
the comparison inside the HSA state model. (ii) The theory is one-sided ($\delta>0$);
the two-sided $\mathrm{BF}_{10}$ is conservative. (iii) A Bayes factor against a point
null is proportional to the prior width by construction, so these numbers are statements
about $\delta=0$ under the stated $\mathcal N(0,0.02^2)$ prior and are not prior-free.

\begin{table}[H]
\centering
\caption{Constant vs.\ time-varying slope: is $\delta\neq0$? Posterior mean of $\delta$
$[95\%$ CI$]$ with the Savage--Dickey Bayes factor $\mathrm{BF}_{10}$ for $\delta\neq0$
within the HSA state model in parentheses (HSA steady, baseline priors,
\textbf{mixed-frequency} design, full sample, unrestricted).
$\mathrm{BF}_{10}>1$ favors a concentration-dependent slope.}
\label{tab:mc}
\resizebox{\linewidth}{!}{\input{../results/tables/annual_q4/delta_grid.tex}}
\end{table}

\paragraph{Verdict.} For the negative unemployment gap, allowing the slope to vary with the firm
count improves the in-sample fit score and the Bayes factor favors $\delta\neq0$
strongly under both CPI indices. For the headline output gaps the gain is marginal,
consistent with the appendix tables: a positive competition dependence appears only in
the core-CPI HP specification, not in BN. For the labor-share gap and inverse markup
there is no inflation signal, so constant and time-varying slopes fit equally well. We
stress that this is an in-sample fit comparison, not evidence of out-of-sample
predictive gains.

\paragraph{Limits of the Savage--Dickey ratios.} We read them only as
the point null $\delta=0$ \emph{within the HSA state model}; they are not a CES-vs-HSA
marginal-likelihood ratio. The in-sample LPPD and RMSE conditioned on the latent states
are likewise not out-of-sample predictions, so they are not used as grounds to select a
more complex model. The model-versus-model question is answered separately, below.

\subsection{Conditional marginal likelihood: CES versus HSA steady}
\label{sec:cml}
The Savage--Dickey ratio asks whether $\delta=0$ inside one model; the fit scores ask
which model tracks this sample better. Neither is a model comparison in the Bayesian
sense. That requires the two models to be densities of the same data given the same
conditioning set,
\begin{equation}
\log p(\pi\mid x,N^{\mathrm{obs}},M)
 = \log m(\pi,N^{\mathrm{obs}},x\mid M) - \log m(N^{\mathrm{obs}}\mid M) - \log m(x),
\label{eq:cml}
\end{equation}
with the latent states integrated out exactly by the Kalman filter and the static
parameters integrated out by Chib's (1995) identity, using the production sampler's own
block conditionals and explicit reduced Gibbs runs.

The two subtracted terms matter. The Gibbs posterior conditions on $\pi$, $N^{\mathrm{obs}}$
\emph{and} $x$: $\sigma_\zeta^2$ is drawn from the activity-equation residuals, and both
$\phi_1$ and $\sigma_\zeta^2$ are ordinate blocks. So the conditioning set is
$(x,N^{\mathrm{obs}})$, and both Occam factors have to be refunded rather than charged.
They separate because $p(N\mid\theta)$ touches only the firm-count block and
$p(x\mid\phi_1,\sigma_\zeta^2)$ only the activity block. $\log m(x)$ is the same number for
both models; we compute it once per cell and debit both with it, rather than cancelling it
by assertion, so that Table~\ref{tab:conditional-ml} shows the two models conditioned on the
identical quantity. CES has no firm-count block and therefore no $\log m(N)$ term.

\input{../results/tables/annual_q4/conditional_ml.tex}

For the main cell the comparison favours HSA steady by
$\Delta\log m=\CmlMainDelta$, a Bayes factor of about \CmlMainBF{}. Across
\CmlMainSeeds{} seeds the Monte Carlo standard error is \CmlMainMcse{}, which moves the
Bayes factor within \CmlMainBFLow{}--\CmlMainBFHigh{} without touching its sign or its
evidence class. That uncertainty is almost entirely the firm-count block: the standard
deviation of $\log m(N)$ across seeds is \CmlSdMN{} against \CmlSdCes{} for the CES term,
the same slow mixing of the AR(2) state that Section~\ref{sec:scope} describes, now
showing up in the reduced runs rather than in the posterior.

The pattern across cells reproduces, by a different route, what $\delta$ already said.
The core-CPI negative-unemployment-gap cell is the strongest; core CPI with the HP output
gap and headline CPI with the unemployment gap are moderate; the BN output gap and all
three PPI cells are indistinguishable from CES. This is an independent check rather than a
restatement: it compares models rather than testing a coefficient, and it is computed from
the marginal likelihood rather than from the posterior of $\delta$.

Two caveats. The estimator is exact only for the linear-Gaussian models: \textbf{HSA full}
is not linear-Gaussian in the joint state, so it has no exact conditional likelihood and
no Gibbs ordinate over its blocks, and the implementation raises rather than returning a
posterior-mean plug-in. And Chib's estimator carries no internal warning when one draw
dominates an ordinate factor; ours refuses to report a factor supported by too few
effective draws, which is how a $660$-log-point error on one seed was caught rather than
published. \CmlNotIdentified{} of the reported cells hit that guard.

\section{What PCHIP interpolation does to the estimates}
\label{sec:pchip}
The conventional alternative to the mixed-frequency design of
Section~\ref{sec:mixed-freq} is to interpolate the annual firm count to a quarterly
series and then use that series \emph{as data}. This section reports that version and
documents what the interpolation costs. Everything else---3 price indices, 3 activity
measures, 5 models, 3 priors, the inverse-markup slope-bias diagnostic, and all MCMC
settings---is identical. CES contains neither $N_t$ nor a latent firm-count state, so its
likelihood and posterior are invariant to the observation scheme and its 16 cells are
shared across both designs.

\subsection{The two designs, side by side}
Setting $\mathcal T_N=\{1,\dots,T\}$ in \eqref{eq:annual-q4-observation}---every quarter
observed---recovers the interpolated filter exactly; the mixed-frequency design is that
same filter with 93 of the 124 firm-count rows removed. The estimator is not the issue.
What differs is what is fed to it.

\begin{table}[H]
\centering
\caption{What the interpolation changes. Posterior quantities for the main cell
(HSA steady, core CPI, negative unemployment gap, baseline priors).}
\label{tab:pchip-vs-mixed}
\small
\input{../results/tables/shared/pchip_vs_mixed.tex}
\end{table}

Three things follow. First, the interpolated series is so smooth that the model drives
$\sigma_N$ down to $\MainSigmaN$ to fit it, which pins $\bar N_t+\hat N_t$ at all 124 quarters
and leaves only the split free. Second, the AR(2) cycle is pushed to a near-unit root,
$\rho_1+\rho_2=\MainRhoSum$, so its mean-reversion term $1-\rho_1-\rho_2$ collapses to
$\MainMeanReversion$.
That term is the \emph{only} thing anchoring the level of $\hat N_t$: adding a constant
$c$ to $\bar N$ and subtracting it from $\hat N$ leaves the firm-count equation, the trend
equation and---through $\kappa_0\to\kappa_0-\delta c$---the inflation equation all exactly
unchanged. When the anchor vanishes the two states become almost perfectly
counter-identified, which is the $\MainStateCorr$ in the table. Across all \RidgeCellCount{}
estimated HSA
cells the correlation between $1-\rho_1-\rho_2$ and corr$(\bar N_0,\hat N_0)$ is
$\RidgeMechanismCorr$,
so this is the mechanism and not a coincidence. Third, the ``cycle'' the interpolated
design recovers has a $\MainArTwoPeriod$-quarter period: it is a second trend competing with
$\bar N_t$
for the same low-frequency variation, not a business-cycle object.

\paragraph{What is \emph{not} affected.} $\delta$ is essentially unchanged---the
interaction $x_t\bar N_t$ that identifies it is insensitive to a level shift in $\bar N$,
which is absorbed by $\kappa_0$. Across the nine HSA-steady baseline cells the two designs
agree on $\delta$ to the third decimal in eight of them. The substantive conclusion of
this paper does not depend on the choice; the reported precision and the interpretation of
the latent states do.

\subsection{Interpolated results, for comparison}
\input{../results/tables/quarterly_interpolated/headline_results.tex}

\begin{table}[H]
\centering
\caption{Interpolated design: model-by-model results for the negative unemployment gap.}
\label{tab:pchip-models}
\tiny
\resizebox{\linewidth}{!}{\input{../results/tables/quarterly_interpolated/unemployment_by_model.tex}}
\end{table}

\begin{figure}[H]
\centering
\includegraphics[width=0.62\linewidth]{quarterly_interpolated/kappa_t_unemployment_by_inflation.png}
\caption{Interpolated design: time-varying slope $\kt$ for the negative unemployment gap.
Compare Figure~\ref{fig:kt_unemp}: the paths are close, but the bands here do not include
the uncertainty from the 93 firm-count values that were interpolated rather than observed.}
\label{fig:kt_unemp_pchip}
\end{figure}

\begin{figure}[H]
\centering
\includegraphics[width=0.88\linewidth]{quarterly_interpolated/competition_decomposition_hsa_steady_core_unemployment.png}
\caption{Interpolated design: firm-count state decomposition. The cycle $\hat N_t$ reaches
$\pm3.4$ ten-log-point units---a $\pm34$ log-point ``cyclical'' swing in the number of
listed firms---which is the level ridge, not an economic cycle. Compare
Figure~\ref{fig:main-state-decomposition}.}
\label{fig:pchip-state-decomposition}
\end{figure}

\subsection{Prior sensitivity under interpolation}
\begin{table}[H]
\centering
\caption{Interpolated design: three-prior sensitivity for the negative unemployment gap.
$\dagger$ marks cells outside the convergence criteria.}
\label{tab:pchip-priors}
\tiny
\resizebox{\linewidth}{!}{\input{../results/tables/quarterly_interpolated/prior_sensitivity_compact.tex}}
\end{table}

\begin{table}[H]
\centering
\caption{Interpolated design: model-by-model convergence tally.}
\label{tab:pchip-convergence}
\small
\input{../results/tables/quarterly_interpolated/convergence_summary.tex}
\end{table}

\begin{table}[H]
\centering
\caption{Interpolated design: scalar / state-path / derived-path diagnostics by model.
Compare Table~\ref{tab:group-convergence}: the worst-mixing scalar here is the trend
drift $n$ in nearly every cell, which is the level ridge in the parameter block.}
\label{tab:group-convergence-pchip}
\tiny
\resizebox{\linewidth}{!}{\input{../results/tables/quarterly_interpolated/group_convergence_diagnostics.tex}}
\end{table}

\paragraph{A prior--data conflict that is specific to interpolation.} The AR(2) prior,
$\rho_1\sim\mathcal N(0.5,\sigma_\rho^2)$ and $\rho_2\sim\mathcal N(-0.5,\sigma_\rho^2)$
truncated to the stationary region, encodes a cycle of about $5.2$ quarters with a
two-quarter half-life, and does so identically under all three prior sets---they differ
only in width. Under the mixed-frequency design the posterior is compatible with it: the
$\rho_1$ posterior sits between $-1.5$ and $+1.6$ prior standard deviations of the prior
mean. Under interpolation it sits $+2.7$ to $+7.1$ prior standard deviations away. The
prior is not obviously wrong; the interpolated series is simply not the object the prior
describes. One consequence is that the tight-prior cells inherit a prior--data tug-of-war
in the AR(2) block, which is where their mixing failures come from.

\section{Scope condition and the limits of identification}
\label{sec:scope}
For honesty we state where the main result does \emph{not} hold and what it can and
cannot claim.
\begin{itemize}
\item \textbf{Not identified for PPI.} Switching to PPI, closer to the theory's
left-hand side, makes $\delta$ include zero for all activity measures. Despite the
measurement limitation that no PPI expectation exists, we cannot claim that the CPI
relationship holds generally at the price-setting stage.
\item \textbf{HSA full does not converge; const-theta now does, in part.} Every
HSA-full cell is outside the acceptance criteria even under Particle Gibbs, so
$\gamma$---which only HSA full estimates---remains undemonstrated. HSA const-theta is a
different case: its earlier failure was a blocking artifact of the alternating FFBS, and
under the exact joint FFBS \ConstThetaPassCount{} of its \ConstThetaCellCount{} cells now
pass the \emph{whole-run} coefficient rule under the interpolated design, against
\AnnualConstThetaPassCount{} under mixed
frequency, where the AR(2) block fails for the reason below. Its $\delta$ agrees in sign
and magnitude with HSA
steady in both. That is a computational fix, not new identification.
\item \textbf{The joint posterior does not converge, though $\delta$ does.} Every
HSA cell in the main design is outside the whole-run acceptance criteria. The failure is
concentrated in the AR(2) block: $(\rho_1,\rho_2)$ is the worst-mixing scalar in
\AnnualArTwoWorstCount{} of
the \AnnualHsaCellCount{} HSA cells, and the latent-state paths fail alongside it. What
does \emph{not} fail
is the coefficient the conclusions rest on. Applied to $\delta$ \emph{alone} rather than to
all \ScalarParameterCount{} scalars at once, the same $\hat R\leq1.01$ / bulk ESS $\geq400$
thresholds are met
in \AnnualSteadyDeltaPassCount{} of the \AnnualSteadyDeltaCellCount{} HSA-steady baseline
cells at a median bulk ESS of \AnnualSteadyDeltaMedianEss{}, and $\delta$ and $\kappa_0$
fail in only \AnnualDeltaFailCount{} and \AnnualKappaZeroFailCount{} of
\AnnualHsaCellCount{} cells here (\DeltaFailCount{} and \KappaZeroFailCount{} under
interpolation).
We therefore report $\delta$ and $\kt$ as results and treat the estimated cycle
$\hat N_t$, its AR(2) coefficients, and anything that depends on them as diagnostic only.
Readers who want the whole-run criterion satisfied should treat every number in this
paper as provisional.
\end{itemize}

\paragraph{The latent-level ridge, and where it comes from.} The firm-count equation
$N^{\mathrm{obs}}_t=\bar N_t+\hat N_t+\nu_t$ constrains the \emph{sum} of the two states
but not the split. Adding a constant $c$ to $\bar N$ and subtracting it from $\hat N$
leaves that equation, the random-walk trend equation, and---through
$\kappa_0\to\kappa_0-\delta c$---the inflation equation all exactly unchanged. The only
thing that resists is the AR(2) cycle equation, which carries no intercept and therefore
pulls $\hat N$ toward mean zero with force $1-\rho_1-\rho_2$; the initial-state prior
$s_0\sim\mathcal N(0,10I_3)$ resists weakly as well.

Earlier revisions of this paper described the resulting ridge as a limitation of the data.
It is better described as a consequence of the interpolated design. Under interpolation
the AR(2) is pushed to a near-unit root and $1-\rho_1-\rho_2$ collapses to
$\MainMeanReversion$, so the
anchor effectively disappears and the posterior correlation between the two states reaches
$\MainStateCorr$. Under the mixed-frequency design used here the same quantity is
$\AnnualMainMeanReversion$ and the
correlation is $\AnnualMainStateCorr$. Across all \RidgeCellCount{} estimated HSA cells the
two quantities correlate at
$\RidgeMechanismCorr$, which identifies the mechanism (Section~\ref{sec:pchip}).

Two things follow. First, $\kappa_0$---the slope at average competition---is identified
only up to the location normalisation of $\bar N$, and under interpolation that
normalisation comes substantially from the initial-state prior rather than from data:
the posterior correlation between the $\bar N$ level and $\kappa_0$ is $\MainLevelKappaCorr$
there
against $\AnnualMainLevelKappaCorr$ here. Statements about the \emph{level} of the
Phillips-curve slope should
be read with that in mind. Second, $\delta$ is not affected, because a level shift in
$\bar N$ is absorbed by $\kappa_0$ and leaves the interaction $x_t\bar N_t$ that identifies
$\delta$ unchanged. That is why the headline result is stable across designs while the
state decomposition is not.

\paragraph{Not a causal effect.} $\bar N_t$ is a near-monotone low-frequency trend, so
on this design $\delta$ cannot be separated from other things that trended over the
same decades: the Volcker monetary-policy regime change, the anchoring of inflation
expectations, changes in labor-market institutions and unionization, globalization and
import competition, changes in the frequency of price adjustment, or a simple time
trend / structural break. The most we can say is that \emph{the low-frequency decline
in the firm count and the decline in inflation's sensitivity to the negative unemployment gap
occurred together, with a sign consistent with the HSA theory}. We therefore avoid
``competition causes flattening'' and ``we identify the mechanism.''

\section{Robustness}
\label{sec:robust}
We judge robustness on three points: (i) how the sign, interval, and magnitude change
under alternative priors, (ii) whether the result survives switching the price index to
PPI, and (iii) whether the same conclusion holds using only converged runs. The prior
check re-estimates \emph{every cell} under baseline, weak, and tight priors
(Table~\ref{tab:current-priors}); it is a genuine re-estimation of the current runs, not
a re-weighting of a single posterior. What the three sets actually change is in
Table~\ref{tab:prior-sets}, and two features of it matter for reading the sweep.
First, the three sets differ in \emph{dispersion only} --- every coefficient prior keeps
its centre across the three --- but the widening is not uniform: it multiplies most
coefficient standard deviations by $2.5$ and $\delta$, $\gamma$ by $5$. (An earlier
revision also moved the centre of $\theta,\theta_0$ from $0.1$ to $0$ between sets, which
mixed a location change into the sweep; all three now centre them at $0$, matching
$\delta$ and $\gamma$ and taking no prior position on the sign the theory predicts.)
Second, ``tight'' holds every variance prior
at the same implied mean while concentrating it, whereas ``weak'' sets the inverse-gamma
shape to $a=1$, for which the prior mean does not exist. That last choice is why the
weak cells, not the tight ones, are where the sampler mixes worst.

\begin{table}[H]
\centering
\caption{What the three prior sets change. Coefficient rows are $N(\text{mean},
\text{sd}^2)$; variance rows are $\operatorname{IG}(a,b)$ with the implied prior mean
$b/(a-1)$ in small type. Read from \code{configs/priors\_*.yaml}, so this table cannot
disagree with the priors the sampler loaded.}
\label{tab:prior-sets}
\small
\resizebox{\linewidth}{!}{\input{../results/tables/shared/prior_sets.tex}}
\end{table}

\begin{table}[H]
\centering
\caption{Three-prior sensitivity under the main mixed-frequency design
(negative-unemployment-gap specification): $\delta$ under baseline, weak, and tight
priors, by model and price index. Each cell is a separate re-estimation, not a
re-weighting. The interpolated counterpart is Table~\ref{tab:pchip-priors}.}
\label{tab:current-priors}
\tiny
\resizebox{\linewidth}{!}{\input{../results/tables/annual_q4/prior_sensitivity_compact.tex}}
\end{table}

\paragraph{HSA steady.} The headline-CPI and core-CPI $\delta$ keep a positive sign
under weak and baseline priors, but not their interval: core CPI moves
$\AnnualCorePriorWeakDelta$ (weak), $\AnnualCoreUnempDelta$ (baseline),
$\AnnualCorePriorTightDelta$ (tight), so the tight prior pulls the interval across zero. The
result is \emph{not} prior-invariant, and the paper does not claim otherwise. PPI includes
zero under all three priors.

Because the three prior sets move every prior at once, we ran a factorial experiment
crossing $\delta$'s own prior with the AR(2) prior while holding all others at baseline,
to see which channel drives the sweep. \textbf{That experiment is run on the interpolated
design}, which is the setting in which the AR(2) prior is most in tension with the data
(Section~\ref{sec:pchip}) and therefore has the most scope to move $\delta$; a smaller
AR(2) contribution under mixed frequency follows a fortiori, but we have not estimated
it. The answer is $\delta$'s own prior: tightening it
alone accounts for $\PriorDecompHeadlineDeltaPriorShare\%$ of the headline-CPI shrinkage
and $\PriorDecompCoreDeltaPriorShare\%$ of the core-CPI
shrinkage, against $\PriorDecompHeadlineArTwoPriorShare\%$ and
$\PriorDecompCoreArTwoPriorShare\%$ for the AR(2) prior (with a
$\PriorDecompCoreInteractionShare\%$ interaction in
the core cell, the two channels partly offsetting). The tension between the AR(2) prior
and the interpolated data is therefore not what is moving $\delta$.

Mixing is a separate matter, and it is the \emph{weak} corner that breaks, not the tight
one. For the main cell the worst scalar has bulk ESS \AnnualWorstEssBaseline{} under
baseline priors and \AnnualWorstEssTight{} under tight, but \AnnualWorstEssWeak{} under
weak, with $\hat R$ rising to \AnnualWorstRhatWeak{}. The weak set puts the
inverse-gamma variance priors at $a=1$ (Table~\ref{tab:prior-sets}), which has no prior
mean; the posterior locations of $\sigma_u,\sigma_\varepsilon$ barely move relative to
baseline, so what the weak set buys is not a different answer but a much heavier tail for
the sampler to explore.

\paragraph{Model variants.} HSA const-theta reproduces the main result almost exactly
under the added entry term: $\delta=\AnnualConstThetaHeadlineDelta$ for headline CPI and
$\delta=\AnnualConstThetaCoreDelta$ for core CPI, against $\AnnualHeadlineUnempDelta$ and
$\AnnualCoreUnempDelta$ under HSA
steady. The agreement is close enough that the firm-count-trend channel is evidently not
being carried by the omission of the entry term. Under the interpolated design the same
comparison also holds, and there \ConstThetaPassCount{} of the \ConstThetaCellCount{}
const-theta cells additionally satisfy the whole-run
coefficient rule; under mixed frequency \AnnualConstThetaPassCount{} do, for the AR(2)
reason above. HSA full
remains outside the criteria in every cell under either design, so $\gamma$ is not
identified here. HSA dynamic is short of ESS in most cells once the trend drift $n$ and
the variance block are included in the rule.

\paragraph{Period robustness.} The state equation assumes consecutive quarters. We do
not use the old treatment that deletes an interior span and compresses what precedes and
follows. The current 77 runs are full-sample comparisons; until a separate contiguous
subsample estimation of the corrected revision is done, we do not include pre-2008 or
post-1988 numbers in the conclusions.

\section{Conclusion}
Over 1982--2012 in the United States, the low-frequency decline in the number of listed
firms and the decline in inflation's sensitivity to the negative unemployment gap moved together.
We read this as theory-consistent time-series evidence, not a causal effect.
\begin{itemize}
\item \textbf{Sign consistent with HSA.} For the core-CPI negative-unemployment-gap specification
$\delta=\AnnualCoreUnempDelta$ ($\mathrm{BF}_{10}=\AnnualCoreUnempDeltaBF$); a falling firm count
implies a flatter slope, as the theory predicts.
\item \textbf{Economically sizable flattening.} The implied $\kt$ falls from
$\AnnualCoreUnempKappaStart$ to $\AnnualCoreUnempKappaEnd$; the start-to-end response to a one-point
negative unemployment gap differs by about $\AnnualCoreUnempKappaDrop$ point.
\item \textbf{Core CPI aligns the output gap.} Under core CPI the HP output gap shows the
same positive sign; BN stays inconclusive.
\item \textbf{Not confirmed for PPI.} The negative-unemployment-gap specification gives
$\delta=\AnnualPPIUnempDelta$, with the interval including zero for all activity measures, so
the relationship is not a general upstream (price-setting-stage) effect.
\item \textbf{No CES-slope bias found.} The theory's predicted negative bias of the CES
slope is not supported: both the CES--HSA difference and the HSA-implied omitted-variable
bias have intervals including zero, with, if anything, positive point estimates.
\item \textbf{Cyclical-entry channel unidentified.} The theory
predicts a positive entry coefficient $\theta$ in $-\theta_t\hat N_t$. In the main design
HSA const-theta gives $\AnnualConstThetaHeadlineTheta$ for headline CPI
($\Pr(\theta>0)=\AnnualConstThetaHeadlineThetaProb$) and
$\AnnualConstThetaCoreTheta$ for core CPI
($\Pr(\theta>0)=\AnnualConstThetaCoreThetaProb$): both centred essentially on zero, with
posteriors that sit on top of the prior in Figure~\ref{fig:pp-unemp}. The data do not
move $\theta$. An earlier revision reported these two cells as carrying \emph{opposite}
signs; that contrast was largely an artefact of a prior centred at the theory's
preferred $+0.1$, and it does not survive centring the prior at zero.
What does survive is a dependence on choices that should not matter: under
interpolation the same core-CPI cell gives $\ConstThetaCoreTheta$ with
$\Pr(\theta>0)=\ConstThetaCoreThetaProb$, against
$\AnnualConstThetaCoreTheta$ here, and within the interpolated design the sign flips
between activity measures ($\ConstThetaCoreTheta$ with the unemployment gap versus
$\ConstThetaCoreBNTheta$ with the BN output gap, the latter putting
$\ConstThetaCoreBNThetaProbNeg\%$ of its mass on the
theory-wrong side). We claim no support for this channel. $\gamma$ is estimated only in
HSA full, which converges in no cell under either design.
\item \textbf{The result does not depend on the interpolation, but the reported
precision does.} Estimating with the annual firm count PCHIP-interpolated to quarterly
and used as data gives $\delta=\CoreUnempDelta$ against $\AnnualCoreUnempDelta$ here---the
same number to within a rounding step. What the interpolation changes is everything around
it: it drives the
firm-count measurement error from $\AnnualMainSigmaN$ to $\MainSigmaN$, forces the AR(2)
cycle to a near-unit root, and through that opens a near-exact ridge between the trend and
cycle states (posterior correlation $\MainStateCorr$ against $\AnnualMainStateCorr$). We report it in
Section~\ref{sec:pchip} as a comparison rather than as the main specification.
\item \textbf{Identification caveat.} The variation is concentrated in the 1980s and
$\bar N_t$ is nearly collinear with a linear trend, so $\delta$ is an association, not a
causal effect.
\end{itemize}

\section{Status and open items}
\label{sec:status}
\paragraph{Estimation status (auto-generated).}
The main results use the \AnnualReportRunCount{} runs of the mixed-frequency design,
of which \AnnualReportWarningCount{} fail the scalar acceptance rule (maximum
$\hat R\leq1.01$ and minimum bulk ESS $\geq400$ over every stored scalar). The
interpolated comparison of Section~\ref{sec:pchip} uses \ReportRunCount{} runs, of which
\ReportWarningCount{} fail the same rule. CES carries no latent firm-count state, so its
16 cells are shared across both designs. All headline estimates, tables, and figures are
generated automatically from the saved posteriors; no numbers from earlier revisions are
read in. Section~\ref{sec:scope} states which failures bear on the conclusions and which
do not.

\paragraph{Settled in the current revision.} The positive $\delta$ and flattening $\kt$
for the CPI negative-unemployment-gap specification of HSA steady, and their insensitivity
to the firm-count observation design; the positive sign of the core-CPI HP specification
and the non-result of BN; the non-results of the three PPI activity specifications; the
prior sensitivity of $\delta$ and its decomposition into $\delta$'s own prior versus the
AR(2) prior; the economic magnitude; the CES--HSA slope-bias diagnostic via the inverse
markup; and the diagnosis that the latent-level ridge is a property of the interpolated
design rather than of the data.

\paragraph{Open (empirical / econometric).}
\begin{itemize}
\item \textbf{Causal identification.} The horse race against a linear time trend cannot
distinguish the two. Before any claim beyond association, add quadratic trends,
structural-break dummies, and placebo interactions.
\item \textbf{Expectation mismatch.} Beyond core-matched expected inflation, connect a
PPI forecast series (e.g.\ the Livingston Survey) so the expectation term matches the
price index.
\item \textbf{Genuinely out-of-sample evaluation.} Table~\ref{tab:fit-comparison} now
reports prequential LPD$_1$, WAIC and PSIS-LOO next to the plug-in score, but all four
are computed on the estimation sample. A rolling out-of-sample exercise---re-estimating
on an expanding window and scoring only genuinely future quarters---remains open, and
would be the only way to settle whether the time-varying slope forecasts better or merely
fits better. The PPI and HSA-full cells additionally have maximum Pareto $\hat k>0.7$, so
their LOO values should be replaced by exact or $K$-fold cross-validation.
\item \textbf{Period robustness.} Re-estimate the current revision using only contiguous
quarters. The pre-2008/post-1988 numbers from earlier revisions are not used for the
conclusions of this paper.
\item \textbf{Bayes-factor robustness.} In addition to the two-sided SDDR, report the
one-sided ($\delta>0$) evaluation and the KDE bandwidth sensitivity.
\item \textbf{Annual-Q4 state block.} The marginal chains of $\delta,\kappa_0$ are good,
but $\rho_1,\rho_2$ and some time points of $\bar N_t,\hat N_t$ have low ESS / high Rhat.
Consider reparameterization, longer independent chains, and priors that relax the weak
identification of the state equation.
\end{itemize}

\paragraph{Open (computational).}
\begin{itemize}
\item \textbf{Conditional marginal likelihood --- now reported
(Section~\ref{sec:cml}), with two things still open.} The implementation integrates the
latent states out exactly by the Kalman filter, uses the estimating model's own initial
state, takes the AR(2) ordinate from the sampler's own initial-lag conditional,
normalizes the truncated $(\rho_1,\rho_2)$ prior and conditional over the stationary
triangle, and factorizes the ordinate with explicit reduced Gibbs runs. What still limits
it: (i) HSA full is excluded by construction --- its bilinear latent term is not
linear-Gaussian, so no exact conditional likelihood or Gibbs ordinate over its blocks
exists, and the implementation raises rather than returning the posterior-\emph{mean}
plug-in an earlier version reported; (ii) the Monte Carlo standard error, \CmlMainMcse{}
on the main cell, is dominated by the firm-count block's reduced runs and would need
longer runs, not a better estimator, to shrink. Only the mixed-frequency design and
baseline priors have been run.
\item \textbf{Weak convergence of full.} HSA full remains outside the criteria in every
cell under Particle Gibbs. HSA const-theta's failure, by contrast, was a blocking
problem and is resolved by the exact joint FFBS; that is a computational fix and is not
evidence about identification.
\item \textbf{Weak identification of $\gamma$.} An identification strategy is needed to
relax the collinearity between $\hat N\bar N$ and $\hat N$.
\item \textbf{Latent-level ridge.} $\bar N_t$ and $\hat N_t$ have posterior correlation
$\MainStateCorr$ under interpolation and $\AnnualMainStateCorr$ under mixed frequency;
only their sum is pinned by the firm-count equation. $\kappa_0$ is
therefore identified partly through the initial-state prior, which is held fixed across
the prior sweep. $\delta$, which loads on the interaction $x_t\bar N_t$, is not affected
by this ridge.
\end{itemize}

\appendix
\section{Appendix: results by output gap (mixed-frequency, main design)}
\label{app:output-gaps}
The body takes the negative unemployment gap as the main activity measure and separates the
output-gap results here. Because HP and BN estimate potential output differently, we do
not combine them into a single ``output gap'' coefficient but report the full baseline
results by price index (3) and model (5) for each definition. The slope is the constant
$\kappa$ for CES/HSA dynamic and the slope at average competition $\kappa_0$ for the
other HSA models. $\delta$ is the sensitivity of the slope to the firm-count trend, and
entry denotes $\theta$ or $\theta_0$.

\begin{figure}[H]
\centering
\includegraphics[width=0.82\linewidth]{annual_q4/delta_output_gaps.png}
\caption{HSA-steady $\delta$ by output gap (posterior mean and $95\%$ CI). HP and BN are
separated and headline CPI, core CPI, and PPI are shown on the same scale.}
\label{fig:output-gap-delta}
\end{figure}

\subsection{HP output gap}
For HSA steady the headline-CPI $\delta$ is near zero, while core CPI is positive:
\begin{equation*}
\delta_{\mathrm{Core,HP}}=\AnnualCoreHPDelta,
\qquad \mathrm{BF}_{10}=\AnnualCoreHPDeltaBF.
\end{equation*}
The PPI $\delta=\AnnualPPIHPDelta$ includes zero. The positive competition dependence of the
HP specification is therefore limited to core CPI and is not a result for price indices
in general.

\begin{table}[H]
\centering
\caption{HP output gap: baseline estimates by price index and model. $\dagger$ marks
failure of the convergence criteria.}
\label{tab:output-gap-hp}
\tiny
\resizebox{\linewidth}{!}{\input{../results/tables/annual_q4/output_gap_hp_by_model.tex}}
\end{table}

\subsection{BN output gap}
For HSA steady the $\delta$ point estimates are negative for headline CPI, positive for
core CPI, and negative for PPI, but all three $95\%$ credible intervals include zero
(headline $\AnnualHeadlineBNDelta$, core $\AnnualCoreBNDelta$, PPI $\AnnualPPIBNDelta$). Each
$\mathrm{BF}_{10}$ is also about one, so the BN specification provides no evidence for a
competition dependence of the slope.

\begin{table}[H]
\centering
\caption{BN output gap: baseline estimates by price index and model. $\dagger$ marks
failure of the convergence criteria.}
\label{tab:output-gap-bn}
\tiny
\resizebox{\linewidth}{!}{\input{../results/tables/annual_q4/output_gap_bn_by_model.tex}}
\end{table}

\paragraph{How to read the whole.} HP core CPI agrees in direction with the CPI result
for the negative unemployment gap, but HP headline, HP PPI, and all price indices of BN do not
support it. And because HSA full has every cell, and const-theta most cells, outside the convergence
criteria, we do not count sign agreement in the complex models as robustness evidence.
The output-gap results are not a claim that ``the result always holds across activity
measures'' but a sensitivity analysis showing which combinations of measurement and
price index the result depends on.

\section{Appendix: results by output gap (interpolated design, for comparison)}
\label{app:annual-output-gaps}
We apply the same mixed-frequency estimation as Section~\ref{sec:main} to
the HP and BN output gaps. Each table gives the full baseline results by price index (3)
and model (5); only the CES row is a common reference that does not depend on the
firm-count observation scheme. We first compare the HSA-steady $\delta$ on the same
scale, then show the full coefficients and diagnostics by model.

\begin{figure}[H]
\centering
\includegraphics[width=0.82\linewidth]{quarterly_interpolated/delta_output_gaps.png}
\caption{Interpolated design: HSA-steady $\delta$ by HP and BN output gap (posterior mean
and $95\%$ CI).}
\label{fig:annual-output-gap-delta}
\end{figure}

\subsection{HP output gap}
For the annual-Q4 HSA steady, headline CPI gives
$\delta=\AnnualHeadlineHPDelta$ ($\mathrm{BF}_{10}=\AnnualHeadlineHPDeltaBF$), core CPI
gives $\delta=\AnnualCoreHPDelta$ ($\mathrm{BF}_{10}=\AnnualCoreHPDeltaBF$), and PPI
gives $\delta=\AnnualPPIHPDelta$ ($\mathrm{BF}_{10}=\AnnualPPIHPDeltaBF$).

\begin{table}[H]
\centering
\caption{Interpolated design: HP output gap, baseline estimates by price index and model.}
\label{tab:annual-output-gap-hp}
\tiny
\resizebox{\linewidth}{!}{\input{../results/tables/quarterly_interpolated/output_gap_hp_by_model.tex}}
\end{table}

\subsection{BN output gap}
For the annual-Q4 HSA steady, headline CPI gives
$\delta=\AnnualHeadlineBNDelta$ ($\mathrm{BF}_{10}=\AnnualHeadlineBNDeltaBF$), core CPI
gives $\delta=\AnnualCoreBNDelta$ ($\mathrm{BF}_{10}=\AnnualCoreBNDeltaBF$), and PPI
gives $\delta=\AnnualPPIBNDelta$ ($\mathrm{BF}_{10}=\AnnualPPIBNDeltaBF$).

\begin{table}[H]
\centering
\caption{Interpolated design: BN output gap, baseline estimates by price index and model.}
\label{tab:annual-output-gap-bn}
\tiny
\resizebox{\linewidth}{!}{\input{../results/tables/quarterly_interpolated/output_gap_bn_by_model.tex}}
\end{table}

These correspond one-to-one with Appendix~\ref{app:output-gaps} of the PCHIP version. In
the comparison we read not only the point estimate but the $95\%$ credible interval, the
Bayes factor, and the convergence flag together. Where an interval widens under
annual-Q4, that is the result of not treating the quarterly firm count as observed and
propagating the interpolation uncertainty.

\section{Appendix: full model-by-model tables (negative unemployment gap)}
\label{app:model-tables}
This appendix collects the dense, auto-generated tables that support
Sections~\ref{sec:bench} and~\ref{sec:robust}: the full model-by-model coefficients for
the negative unemployment gap and the whole-run convergence tally. (The three-prior sensitivity
table now sits with the robustness discussion, Table~\ref{tab:current-priors}.)
The slope is $\kappa$ for CES/HSA dynamic and $\kappa_0$ for the other models; entry is
$\theta$ or $\theta_0$; $\dagger$ marks cells outside the convergence criteria.

\begin{table}[H]
\centering
\caption{Model-by-model results for the negative-unemployment-gap specification. Slope is
$\kappa$ for CES/HSA dynamic and $\kappa_0$ otherwise; entry is $\theta$ or $\theta_0$.
$\dagger$ marks cells that fail the convergence criteria and are not interpreted.}
\label{tab:current-models}
\tiny
\resizebox{\linewidth}{!}{\input{../results/tables/annual_q4/unemployment_by_model.tex}}
\end{table}

\begin{figure}[H]
\centering
\includegraphics[height=0.68\textheight,keepaspectratio]{annual_q4/slope_by_model_inflation.png}
\caption{Slope at average competition for the negative-unemployment-gap specification. $\kappa$
for CES/HSA dynamic and $\kappa_0$ otherwise. Circles are runs meeting the convergence
criteria; gray crosses are runs to watch.}
\label{fig:current-slope-model}
\end{figure}

\begin{figure}[H]
\centering
\includegraphics[width=0.94\linewidth]{annual_q4/alpha_entry_by_model_inflation.png}
\caption{Inflation inertia $\alpha$ (left) and entry coefficient $\theta$ or $\theta_0$
(right) for the negative-unemployment-gap specification. Circles and gray crosses are as in
Figure~\ref{fig:current-slope-model}.}
\label{fig:alpha-entry-main}
\end{figure}

\begin{table}[H]
\centering
\caption{Model-by-model convergence tally across all \ReportRunCount{} runs.
``warnings'' counts runs failing the \emph{scalar} rule, ``state warnings'' the
latent-path rule, and ``joint warnings'' at least one of the three groups
(Section~\ref{sec:model}). The state-block column names the sampler that drew
$\bar N_t,\hat N_t$.}
\label{tab:main-convergence}
\small
\input{../results/tables/annual_q4/convergence_summary.tex}
\end{table}

\begin{table}[H]
\centering
\caption{Group-by-group convergence diagnostics for the negative-unemployment-gap
baseline cells: scalar parameters, the trend drift $n$ on its own, the latent-state
paths $\bar N_t,\hat N_t$, and the derived paths $\kt,\theta_t$. ``worst scalar'' names
the scalar with the lowest bulk ESS. Thresholds are $\hat R\leq1.01$ and bulk
ESS $\geq400$. Mixed-frequency (main) design. This table is the evidence behind the
$\dagger$ and ``OK (coef)'' marks elsewhere: it shows whether a failure is driven by a
coefficient of interest or by a nuisance term. Here the worst-mixing scalar is the AR(2)
block in almost every cell, which for HSA steady is a nuisance object.}
\label{tab:group-convergence}
\tiny
\resizebox{\linewidth}{!}{\input{../results/tables/annual_q4/group_convergence_diagnostics.tex}}
\end{table}

\section{Appendix: every run behind this paper}
\label{app:runs}
The tables above report cells; this appendix reports \emph{runs}. Table~\ref{tab:run-manifest-annual}
and Table~\ref{tab:run-manifest-pchip} list every posterior the paper reads --- one row
per (model, data spec, prior) cell, with the run identifier, the sampler that drew the
state block, the estimation sample length, and the three convergence groups. Every number
elsewhere in this paper is computed from exactly these runs and no others.

The run directory holds exactly these runs and nothing else: one directory per
(model, data spec, prior, observation design), named for that cell and carrying no
timestamp, so the directory listing and the tables below are the same set. Superseded
estimation revisions are deleted rather than accumulated; the revision string in each
run's metadata records which vintage of the inputs produced it.

\begin{table}[H]
\centering
\caption{Mixed-frequency (main) design: every run the paper reads. ``state sampler'' is
the sampler that drew $\bar N_t,\hat N_t$; ``max Rhat''/``min bulk ESS'' are over the
\ScalarParameterCount{} scalars, and the state columns are over the latent paths.}
\label{tab:run-manifest-annual}
\tiny
\resizebox{\linewidth}{!}{\input{../results/tables/annual_q4/report_run_manifest.tex}}
\end{table}

\begin{table}[H]
\centering
\caption{Interpolated (PCHIP) design: every run the paper reads. CES carries no latent
firm-count state, so its cells are shared with the mixed-frequency design rather than
re-estimated.}
\label{tab:run-manifest-pchip}
\tiny
\resizebox{\linewidth}{!}{\input{../results/tables/quarterly_interpolated/report_run_manifest.tex}}
\end{table}

\section{Appendix: prior against posterior, every model and coefficient}
\label{app:prior-posterior}
The tables report intervals; these figures report the whole marginal posterior, and
put it next to the prior it was estimated under. Each figure covers one activity
measure under the main mixed-frequency design: models down the rows, the four
coefficients across the columns, and the three price indices as the three curves in
each panel. A dash marks a coefficient the model restricts away. The dashed line is
the prior, read from each run's own saved prior block rather than from the config
file, so a curve and its prior are always the pair that actually produced the draws.

Three things are worth reading off them. First, down the $\delta$ column, HSA steady,
HSA const-theta and HSA full give an almost identical posterior: adding the entry
term does not move the coefficient the conclusions rest on. Second, across every
panel, the core-CPI curve is the sharpest and the PPI curve the widest --- the same
ordering the intervals show, but here it is visibly a difference in how much the data
say rather than a difference in point estimates. Third, and most directly, the
$\gamma$ column sits on top of its own prior in all three price indices. The
posterior is the prior; the data are silent about $\gamma$. That is the same
conclusion Section~\ref{sec:scope} states from the convergence diagnostics, arrived
at without reference to mixing.

\begin{figure}[H]
\centering
\includegraphics[width=0.84\linewidth]{annual_q4/prior_vs_posterior_unemployment_gap.png}
\caption{Prior against posterior, negative unemployment gap (the main activity
measure), mixed-frequency design, baseline priors.}
\label{fig:pp-unemp}
\end{figure}

\begin{figure}[H]
\centering
\includegraphics[width=0.84\linewidth]{annual_q4/prior_vs_posterior_output_gap_hp.png}
\caption{Prior against posterior, HP output gap, mixed-frequency design, baseline
priors.}
\label{fig:pp-hp}
\end{figure}

\begin{figure}[H]
\centering
\includegraphics[width=0.84\linewidth]{annual_q4/prior_vs_posterior_output_gap_bn.png}
\caption{Prior against posterior, BN output gap, mixed-frequency design, baseline
priors. The interpolated-design counterparts of all three figures are written to
\code{results/figures/quarterly\_interpolated/} by the same builder.}
\label{fig:pp-bn}
\end{figure}

\end{document}
